\section{Experiments}
\label{sec:experiments}

Evaluating repository-to-tool conversion against \emph{unknown} objectives is
hard: the richest signal about what a repository can do comes from its authors.
We therefore adopt \textbf{TM-Bench}, released with ToolMaker
\citep{toolmaker2025}, one of the few tool-/MCP-creation benchmarks with
objectives known to be achievable and human-authored tests for the tools
themselves, a repository's own test suite would not cover tool usage beyond
what the authors happened to test. Its fifteen tasks were curated with medical
and life-science researchers to reflect realistic problems, concentrated in
pathology, radiology, and omics, and rounded out with 3D vision, imaging,
tabular analysis, and NLP for broader coverage.

\paragraph{Setup.} Alembic runs on \texttt{openrouter/z-ai/glm-5.2} in
\emph{known-target} mode (it still discovers every workflow it can, but is also
told to expose a tool matching each TM-Bench task) and exports self-contained
artefacts (\texttt{code.py} and the recorded \texttt{setup.sh}) scored inside
ToolMaker's harness on the same held-out pytest suite. For the ToolMaker
baseline we quote the published GPT-4o figures; a matched GLM-5.2 re-run did not
yet reproduce them, so, as a system demonstration, we report those numbers for
context and leave a controlled same-model comparison to future work. The
reference harness needed only peripheral fixes (download-command flags,
container leakage on failure, importing Alembic's Python-3.11 artefacts into the
harness's Python-3.12 scoring environment, cost-untracked models); our fork:
\url{https://github.com/stas1f1/ToolMaker}.

\paragraph{Data.} Two STAMP tasks run over TCGA whole-slide cohorts whose
reference CSVs list $\approx$1.5\,TB of slides, though the tests need only ten
CRC and two named slides; since that subset cannot be guessed beforehand, we
follow the TM-Bench authors in disabling dataset downloads here while still
allowing weight downloads and in-container cloning of the task data. Both tasks
ride a \emph{single} Alembic run (the Explorer scopes both tools, the Coder
implements both), falling back to separate runs only on failure.

\section{Results and Discussion}
\label{sec:results}

% Auto-generated fragment — \input this into the paper; do not edit by hand.
\begin{table}[t]
\centering
\footnotesize
\setlength{\tabcolsep}{4.5pt}
\renewcommand{\arraystretch}{1.15}
\resizebox{\columnwidth}{!}{%
\begin{tabular}{@{}ll rrr rrr@{}}
\toprule
& & \multicolumn{3}{c}{\textbf{alembic (ours)}} & \multicolumn{3}{c}{\textbf{ToolMaker} \small(W\"olflein et al., 2025)} \\
\cmidrule(lr){3-5}\cmidrule(lr){6-8}
& Task & Invoc. & Tests & Tokens & Invoc. & Tests & Tokens \\
\midrule
\multirow{6}{*}{\rotatebox[origin=c]{90}{Pathology}} & conch\_extract\_features & \G{3/3} & \G{9/9} & 1,206,593 & \G{3/3} & \G{9/9} & 171,226 \\
 & musk\_extract\_features & \G{3/3} & \G{6/6} & 1,098,224 & \G{3/3} & \G{6/6} & 696,386 \\
 & pathfinder\_verify\_biomarker & \R{0/2} & \Y{4/6} & 956,298 & \R{0/2} & \Y{4/6} & 356,825 \\
 & stamp\_extract\_features & \G{3/3} & \G{12/12} & 1,758,210$^\dagger$ & \G{3/3} & \G{12/12} & 631,138 \\
 & stamp\_train\_classification\_model & \R{0/3} & \R{0/9} & 1,758,210$^\dagger$ & \G{3/3} & \G{9/9} & 1,249,521 \\
 & uni\_extract\_features & \G{3/3} & \G{9/9} & 1,683,079 & \G{3/3} & \G{9/9} & 326,806 \\
\midrule
\multirow{2}{*}{\rotatebox[origin=c]{90}{Rad.}} & medsam\_inference & \G{3/3} & \G{6/6} & 3,442,100 & \G{3/3} & \G{6/6} & 508,954 \\
 & nnunet\_train\_model & \R{0/2} & \Y{2/4} & 5,350,444 & \R{0/2} & \R{0/4} & 1,792,291 \\
\midrule
\multirow{2}{*}{\rotatebox[origin=c]{90}{Omics}} & cytopus\_db & \G{3/3} & \G{12/12} & 1,967,988 & \G{3/3} & \G{12/12} & 185,912 \\
 & esm\_fold\_predict & \G{3/3} & \G{15/15} & 8,735,809 & \Y{2/3} & \Y{13/15} & 336,754 \\
\midrule
\multirow{5}{*}{\rotatebox[origin=c]{90}{Other}} & flowmap\_overfit\_scene & \R{0/2} & \R{0/6} & 1,704,245 & \G{2/2} & \G{6/6} & 358,552 \\
 & medsss\_generate & \R{0/3} & \R{0/6} & 3,970,326 & \G{3/3} & \G{6/6} & 282,771 \\
 & modernbert\_predict\_masked & \G{3/3} & \G{9/9} & 2,909,347 & \G{3/3} & \G{9/9} & 356,228 \\
 & retfound\_feature\_vector & \R{0/3} & \R{0/6} & 721,849 & \G{3/3} & \G{6/6} & 561,936 \\
 & tabpfn\_predict & \G{3/3} & \G{9/9} & 1,550,724 & \G{3/3} & \G{9/9} & 95,257 \\
\midrule
 & \textbf{Total} & \textbf{27/42} & \textbf{93/124} & \textbf{37,055,236} & \textbf{37/42} & \textbf{116/124} & \textbf{7,910,557} \\
\bottomrule
\end{tabular}}
\caption{TM-Bench-validated comparison against the published \textsc{ToolMaker} baseline, per task, scored on the shared pytest suite; green/yellow/red mark all/some/none passing. $^\dagger$STAMP's two tasks share one Alembic run (tokens counted once).}
\label{tab:alembic-tmbench-compare}
\end{table}

% Auto-generated fragment — \input this into the paper; do not edit by hand.

\begin{table}[t]
\centering
\small
\begin{tabular}{@{}lr@{}}
\toprule
Error class & Count \\
\midrule
Other & 8 \\
FileNotFound & 7 \\
AttributeError & 4 \\
ValueError & 4 \\
KeyError & 3 \\
ModuleNotFound & 2 \\
Import & 1 \\
Runtime & 1 \\
\midrule
\textbf{Total} & \textbf{30} \\
\bottomrule
\end{tabular}
\caption{Runtime exceptions by class across the 14 repositories (30 total).}
\label{tab:alembic-failures}
\end{table}

% Auto-generated fragment — \input this into the paper; do not edit by hand.

\begin{table}[t]
\centering
\small
\begin{tabular}{@{}lrrr@{}}
\toprule
Reason & Env. & Coder & Total \\
\midrule
max\_steps & 2 & 3 & 5 \\
tool\_cycle & 4 & 0 & 4 \\
tool\_repeat & 0 & 3 & 3 \\
\midrule
\textbf{Total} & \textbf{6} & \textbf{6} & \textbf{12} \\
\bottomrule
\end{tabular}
\caption{Cycle-guard aborts by stage (12 total): the reason each stage's inner agent loop was force-stopped, recorded once per stage, distinct from the stage-gate retries, which count separately.}
\label{tab:alembic-aborts}
\end{table}

% Auto-generated fragment — \input this into the paper; do not edit by hand.
\begin{table*}[t]
\centering
\footnotesize
\setlength{\tabcolsep}{4.5pt}
\renewcommand{\arraystretch}{1.15}
\resizebox{\textwidth}{!}{%
\begin{tabular}{@{}ll rrr rrr c rrr@{}}
\toprule
& & \multicolumn{3}{c}{\textbf{Tools}} & \multicolumn{3}{c}{\textbf{Checks}} & & & & \\
\cmidrule(lr){3-5}\cmidrule(lr){6-8}
& Repo & Stable & \makecell{Pass\\Invoc.} & Total & Tests & Exec & Invoc & \makecell{All\\stages?} & Time & \makecell{Stage\\retries} & Tokens \\
\midrule
\multirow{5}{*}{\rotatebox[origin=c]{90}{Pathology}} & CONCH \small\citep{lu2024conch} & \G{5} & \Y{4} & 5 & \G{16/16} & \G{4/4} & \G{4/4} & \G{\cmark} & 24m27s & 0 & 1,206,593 \\
 & MUSK \small\citep{xiang2025musk} & \G{5} & 0 & 5 & \G{10/10} & \G{5/5} & -- & \G{\cmark} & 20m54s & 0 & 1,098,224 \\
 & PathFinder \small\citep{liang2023pathfinder} & \G{4} & \Y{3} & 4 & \G{14/14} & \G{4/4} & \G{3/3} & \G{\cmark} & 26m08s & 2 & 956,298 \\
 & STAMP \small\citep{elnahhas2024stamp} & \Y{4} & \R{0} & 6 & \G{13/13} & \Y{4/6} & \R{0/4} & \G{\cmark} & 56m23s & 0 & 1,758,210 \\
 & UNI \small\citep{chen2024uni} & \G{5} & \G{5} & 5 & \G{10/10} & \G{5/5} & \G{5/5} & \G{\cmark} & 47m46s & 0 & 1,683,079 \\
\midrule
\multirow{2}{*}{\rotatebox[origin=c]{90}{Rad.}} & MedSAM \small\citep{ma2024medsam} & \G{5} & \Y{4} & 5 & \G{10/10} & \G{3/3} & \G{9/9} & \G{\cmark} & 57m18s & 2 & 3,442,100 \\
 & nnU-Net \small\citep{isensee2020nnunet} & \Y{5} & \R{0} & 8 & \G{23/23} & \Y{5/8} & \R{0/2} & \G{\cmark} & 1h12m & 1 & 5,350,444 \\
\midrule
\multirow{2}{*}{\rotatebox[origin=c]{90}{Omics}} & Cytopus \small\citep{kunes2023cytopus} & \G{9} & \G{9} & 9 & \G{24/24} & \G{9/9} & \G{19/19} & \G{\cmark} & 22m20s & 0 & 1,967,988 \\
 & ESM \small\citep{verkuil2022esm1,hie2022esm2} & \G{7} & \Y{4} & 7 & \G{15/15} & \G{7/7} & \G{4/4} & \G{\cmark} & 2h45m & 4 & 8,735,809 \\
\midrule
\multirow{5}{*}{\rotatebox[origin=c]{90}{Other}} & FlowMap \small\citep{smith2024flowmap} & \Y{3} & \Y{1} & 5 & \G{15/15} & \Y{3/5} & \Y{3/4} & \G{\cmark} & 33m52s & 0 & 1,704,245 \\
 & MedSSS \small\citep{jiang2025medsss} & \Y{4} & \Y{1} & 6 & \G{12/12} & \R{0/2} & \Y{2/6} & \G{\cmark} & 1h03m & 0 & 3,970,326 \\
 & ModernBERT \small\citep{warner2024modernbert} & \Y{4} & \R{0} & 5 & \G{14/14} & \R{0/1} & \R{0/2} & \G{\cmark} & 53m27s & 1 & 2,909,347 \\
 & RETFound \small\citep{zhou2023retfound} & \G{3} & 0 & 3 & \G{8/8} & \G{3/3} & -- & \G{\cmark} & 55m38s & 0 & 721,849 \\
 & TabPFN \small\citep{hollmann2025tabpfn} & \G{4} & \G{4} & 4 & \G{13/13} & \G{1/1} & \G{8/8} & \Y{soft-conf.} & 35m33s & 1 & 1,550,724 \\
\midrule
 & \textbf{Total} (14 repos) & \textbf{67} & \textbf{35} & \textbf{77} & \textbf{197/197} & \textbf{53/63} & \textbf{57/70} & \textbf{13\,\cmark}\,+\,1\,soft & \textbf{12h15m} & \textbf{11} & \textbf{37,055,236} \\
\bottomrule
\end{tabular}}
\caption{Alembic self-metrics on the TM-Bench repositories, one row per repository; green/yellow/red mark all/some/none passing. \emph{All stages?} marks an end-to-end delivery (\cmark); TabPFN's environment resolved only after a soft dependency conflict (\emph{soft-conf.}).}
\label{tab:alembic-selfmetrics}
\end{table*}


Table~\ref{tab:alembic-tmbench-compare} reports the TM-Bench-validated
comparison, per task, against the ToolMaker baseline. \emph{Invoc.} is the
fraction of invocation examples whose every unit test passes and \emph{Tests}
the fraction of unit tests that pass, both on the identical TM-Bench pytest
suite, so the denominators match ToolMaker's. Alembic matches or exceeds the
baseline on the pathology and omics tasks and wins outright on \texttt{esm} and
\texttt{nnunet}, while the external suite also exposes tasks that pass Alembic's
own validation but not TM-Bench's held-out checks.

\paragraph{Failure analysis.} Tables~\ref{tab:alembic-failures}
and~\ref{tab:alembic-aborts} break the run down by error type and by cycle-guard
aborts (an agent loop force-stopped for exhausting its step budget or repeating a
tool call). Runtime failures are dominated by missing files and attribute/value
errors rather than by hallucinated symbols (the class the plan gate eliminates
by construction), and the guard aborts fall entirely in the environment and
coder stages, where dependency resolution is hardest. Tasks split in two: clean
or one-slip tasks (\texttt{cytopus}, \texttt{conch}, \texttt{uni}, \texttt{esm},
\texttt{tabpfn}) resolve in a single pass, while deep-cascade tasks hit a chain of
Python-3.11$\to$3.12 dependency conflicts, the failure mode of transcribing a
dependency set that only resolved under the repository's original interpreter,
which is why Alembic preserves that interpreter in a split environment.

Table~\ref{tab:alembic-selfmetrics} reports Alembic's own, richer metrics, one
row per repository. It splits the tools each run emits into \emph{stable} (builds,
imports, passes unit tests) and \emph{perfect} (also passes live invocation),
records passed/attempted \emph{tests}, executions (\emph{exec}) and invocation
checks (\emph{invoc}), and adds wall-clock time, guard-triggered stage retries,
and total token usage. Alembic delivers a served MCP for all fourteen
repositories; only TabPFN needs a soft dependency-conflict resolution.

\subsection{Case study: MedSAM}
\label{sec:casestudy}

MedSAM \citep{ma2024medsam} adapts the Segment Anything model to medical images.
From the URL alone, Alembic's Explorer discovered five tools (inference,
training, preprocessing, scoring, conversion), all passing and four perfectly.
Consider the inference tool (Figure~\ref{fig:casestudy}).

\begin{figure}[!ht]
  \centering
  \includegraphics[width=0.82\columnwidth]{medsam_casestudy.pdf}
  \caption{The synthesized \texttt{medsam\_inference} on the repository's own
  \texttt{img\_demo.png}: scan and box are shelled into the repo environment
  (ViT-B, CPU); the mask shown is the real MedSAM output.}
  \label{fig:casestudy}
\end{figure}

The Explorer grounded it in the repository's ``Get Started'' walkthrough
(\texttt{MedSAM\_Inference.py} on the bundled \texttt{img\_demo.png}, writing a
\texttt{uint8} mask), and the plan gate confirmed the script and its sample
arguments. The Coder wrote a thin
\texttt{def medsam\_inference(...) -> dict} that shells into the repo environment
(headless Matplotlib, CPU, ViT-B); as the mask is a side-effect, it returns
\texttt{\{\}}. Validation invoked it on the demo and a held-out PNG, got
expected-shape masks, and marked it \emph{perfect} (3/3 invocations, 6/6 tests,
no debugger rounds): a published inference script rendered a validated,
callable MCP tool.
